%

\pol{\subsection{Symetrie osiowe}}
\ita{\subsection{Simmetrie assiali}}

\begin{definition}[\pol{symetria osiowa}\ita{simmetria assiale}]
    \pol{Niech $k$ będzie ustaloną prostą, $X$ punktem, zaś $X'$ jego rzutem na prostą $k$.}
    \ita{Sia $k$ una retta fissata, $X$ un punto e $X'$ la sua proiezione sulla retta $k$.}
    \pol{Oznaczmy przez $w$ wektor $2X X'$.}
    \ita{Indichiamo con $w$ il vettore $2X X'$.}
    \pol{Wtedy funkcję $S_k$ zadaną przez}
    \ita{Allora la funzione $S_k$ definita da}
    \begin{equation}
        S_k(X) = T_{w}(X)
    \end{equation}
    \pol{nazywamy symetrią osiową o osi $k$.}
    \ita{si chiama simmetria assiale avente asse $k$.}
\end{definition}

\pol{(Chodzi o to, by rzut punktu $X$ na prostą $k$ był środkiem odcinka łączącego argument oraz wartość funkcji $S_k$).}
\ita{(Vale a dire che la proiezione di $X$ sulla retta $k$ è il punto medio del segmento che unisce l'argomento e il valore della funzione $S_k$).}
\pol{Symetrie osiowe nie mają estetycznego wzoru, jeśli nie chcemy wprowadzać do książki macierzy: punkt $(x_0, y_0)$ odbity względem osi o równaniu $ax + by + c = 0$ przechodzi na}
\ita{Le simmetrie assiali non ammettono una formula particolarmente elegante, se non si vogliono introdurre le matrici: il punto $(x_0, y_0)$ riflesso rispetto alla retta di equazione $ax + by + c = 0$ viene mandato in}
\begin{equation}
    \left(
    x_0 - 2a \cdot \frac{ax_0 + by_0 + c}{a^2 + b^2},
    y_0 - 2b \cdot \frac{ax_0 + by_0 + c}{a^2 + b^2}
    \right),
\end{equation}

\begin{proposition}
\label{kordos_banach_6}%
    \pol{Złożenie dwóch symetrii osiowych, których osie są równoległe, jest translacją (o wektor dwa razy dłuższy niż odległość między osiami, prostopadły do obydwu).}
    \ita{La composizione di due simmetrie assiali con assi paralleli è una traslazione (di vettore perpendicolare a entrambi gli assi e di lunghezza doppia della distanza tra essi).}
\end{proposition}

\begin{proposition}
\label{klkm1}%
    \pol{Niech $k, l$ będą prostymi.}
    \ita{Siano $k$ e $l$ due rette.}
    \pol{Wtedy $S_k \circ S_l \circ S_k = S_m$, gdzie $m$ to prosta będąca odbiciem prostej $l$ względem $k$.}
    \ita{Allora $S_k \circ S_l \circ S_k = S_m$, dove $m$ è la retta ottenuta riflettendo la retta $l$ rispetto a $k$.}
\end{proposition}

\pol{Jeśli figura ma dokładnie dwie osi symetri, to muszą być prostopadłe.}
\ita{Se una figura possiede esattamente due assi di simmetria, essi devono essere perpendicolari.}

%